% ============================================================
%  !Blood – Presentación Beamer
%  Compilar con: pdfLaTeX
%  Subir a Overleaf como nuevo proyecto
% ============================================================
\documentclass[aspectratio=169, 12pt]{beamer}

% ── Codificación & Idioma ────────────────────────────────
\usepackage[utf8]{inputenc}
\usepackage[T1]{fontenc}
\usepackage[spanish]{babel}

% ── Fuentes ──────────────────────────────────────────────
\usepackage{lmodern}

% ── Paquetes visuales ────────────────────────────────────
\usepackage{xcolor}
\usepackage{tikz}
\usetikzlibrary{shapes.geometric, arrows.meta, positioning, calc, shadows}
\usepackage{graphicx}
\usepackage{booktabs}
\usepackage{array}
\usepackage{colortbl}
\usepackage{multicol}
\usepackage{tcolorbox}
\tcbuselibrary{skins, breakable}
\usepackage{fontawesome5}
\usepackage{enumitem}
\usepackage{caption}

% ════════════════════════════════════════════════════════
%  PALETA DE COLORES
% ════════════════════════════════════════════════════════
\definecolor{bloodred}{HTML}{E53935}
\definecolor{blooddk}{HTML}{B71C1C}
\definecolor{bloodlt}{HTML}{FF6659}
\definecolor{bgdark}{HTML}{0D0D0D}
\definecolor{bgcard}{HTML}{1C1C1C}
\definecolor{warmgray}{HTML}{AAAAAA}
\definecolor{offwhite}{HTML}{F5F5F5}
\definecolor{accentpink}{HTML}{FF8A80}

% ════════════════════════════════════════════════════════
%  TEMA BEAMER PERSONALIZADO
% ════════════════════════════════════════════════════════
\usetheme{default}
\usecolortheme{default}
\usefonttheme{professionalfonts}
\setbeamertemplate{navigation symbols}{}
\setbeamertemplate{itemize item}{\color{bloodred}\small\faCircle}
\setbeamertemplate{itemize subitem}{\color{bloodlt}\scriptsize\faCircle}

% ── Colores del tema ─────────────────────────────────────
\setbeamercolor{background canvas}{bg=bgdark}
\setbeamercolor{normal text}{fg=offwhite}
\setbeamercolor{alerted text}{fg=bloodlt}
\setbeamercolor{structure}{fg=bloodred}
\setbeamercolor{block title}{fg=offwhite, bg=blooddk}
\setbeamercolor{block body}{fg=offwhite, bg=bgcard}
\setbeamercolor{block title alerted}{fg=offwhite, bg=bloodred}
\setbeamercolor{block body alerted}{fg=offwhite, bg=bgcard}
\setbeamercolor{frametitle}{fg=offwhite, bg=bgdark}
\setbeamercolor{footline}{fg=warmgray}
\setbeamercolor{title}{fg=offwhite}
\setbeamercolor{subtitle}{fg=warmgray}
\setbeamercolor{author}{fg=offwhite}
\setbeamercolor{institute}{fg=warmgray}
\setbeamercolor{date}{fg=warmgray}
\setbeamercolor{separation line}{bg=bloodred}

% ── Frametitle ───────────────────────────────────────────
\setbeamertemplate{frametitle}{%
  \nointerlineskip
  \begin{beamercolorbox}[wd=\paperwidth, ht=2.6ex, dp=1.4ex, leftskip=12pt]{frametitle}%
    \usebeamerfont{frametitle}\insertframetitle%
    \ifx\insertframesubtitle\@empty\else%
      \hspace{6pt}\usebeamerfont{framesubtitle}{\color{warmgray}\small|\space\insertframesubtitle}%
    \fi%
  \end{beamercolorbox}%
  \begin{beamercolorbox}[wd=\paperwidth, ht=0.4pt]{separation line}\end{beamercolorbox}%
  \vspace{2pt}%
}

% ── Pie de página ────────────────────────────────────────
\setbeamertemplate{footline}{%
  \begin{beamercolorbox}[wd=\paperwidth, ht=2.2ex, dp=1.2ex]{footline}%
    \hspace{12pt}%
    {\color{bloodred}\tiny\faDroplet}\;%
    {\color{warmgray}\tiny !Blood\enspace|\enspace Ana María Campo Agredo}%
    \hfill%
    {\color{warmgray}\tiny\insertframenumber\,/\,\inserttotalframenumber}%
    \hspace{12pt}%
  \end{beamercolorbox}%
}

% ── Fonts ────────────────────────────────────────────────
\setbeamerfont{frametitle}{size=\large, series=\bfseries}
\setbeamerfont{title}{size=\Huge, series=\bfseries}
\setbeamerfont{subtitle}{size=\large, series=\normalfont}
\setbeamerfont{block title}{size=\small, series=\bfseries}
\setbeamerfont{block body}{size=\footnotesize}

% ════════════════════════════════════════════════════════
%  CAJAS PERSONALIZADAS (tcolorbox)
% ════════════════════════════════════════════════════════

% Caja roja (alerta / dato importante)
\tcbset{
  redbox/.style={
    enhanced, colback=bgcard, colframe=bloodred,
    coltext=offwhite, fonttitle=\bfseries\small,
    coltitle=offwhite, attach boxed title to top left={yshift=-2mm, xshift=4mm},
    boxed title style={colback=bloodred, colframe=bloodred, rounded corners},
    rounded corners, arc=6pt,
    left=6pt, right=6pt, top=8pt, bottom=6pt,
    drop shadow={opacity=.3}
  },
  graybox/.style={
    enhanced, colback=bgcard, colframe=bgcard,
    coltext=offwhite, rounded corners, arc=6pt,
    borderline west={3pt}{0pt}{bloodred},
    left=8pt, right=6pt, top=6pt, bottom=6pt,
  },
  statbox/.style={
    enhanced, colback=bgcard, colframe=bloodred!30,
    coltext=offwhite, rounded corners, arc=8pt,
    left=8pt, right=8pt, top=8pt, bottom=8pt,
    halign=center, valign=center,
  }
}

% ── Macros útiles ────────────────────────────────────────
\newcommand{\redtext}[1]{{\color{bloodlt}#1}}
\newcommand{\greytext}[1]{{\color{warmgray}#1}}
\newcommand{\tagbadge}[1]{%
  \tikz[baseline=(t.base)]
  \node[fill=bloodred!15, draw=bloodred!40, rounded corners=10pt,
        inner xsep=6pt, inner ysep=2pt, text=accentpink,
        font=\bfseries\tiny\ttfamily] (t) {\MakeUppercase{#1}};%
}

% ── Gota SVG en TikZ ─────────────────────────────────────
\newcommand{\blooddrop}[1][1.2cm]{%
  \begin{tikzpicture}[baseline=(d.south)]
    \fill[top color=bloodlt, bottom color=blooddk, drop shadow={opacity=.5,shadow xshift=0,shadow yshift=-1pt}]
      (0,#1) .. controls (-0.55*#1,0.35*#1) and (-0.7*#1,-0.3*#1) ..
      (0,-0.7*#1) .. controls (0.7*#1,-0.3*#1) and (0.55*#1,0.35*#1) ..
      (0,#1) coordinate (d);
    \fill[white, opacity=.2]
      (-0.18*#1,0.18*#1) .. controls (-0.35*#1,0.05*#1) and (-0.32*#1,-0.12*#1) ..
      (-0.15*#1,-0.2*#1) .. controls (-0.02*#1,-0.26*#1) and (0.05*#1,-0.18*#1) ..
      (-0.18*#1,0.18*#1);
  \end{tikzpicture}%
}

% ════════════════════════════════════════════════════════
%  METADATOS
% ════════════════════════════════════════════════════════
\title{{\large\faDroplet}\quad !Blood}
\subtitle{Aplicación Móvil de Donación de Sangre\\[4pt]
  \small Informe de Diseño y Especificación Funcional}
\author{Ana María Campo Agredo}
\institute{Universidad del Cauca\\
  Facultad de Ciencias Naturales, Exactas y de la Educación\\
  Departamento de Física --- Popayán, Colombia}
\date{2026}

% ════════════════════════════════════════════════════════
\begin{document}
% ════════════════════════════════════════════════════════

% ──────────────────────────────────────────────────────
%  DIAPOSITIVA 1 – PORTADA
% ──────────────────────────────────────────────────────
{
\setbeamertemplate{footline}{}
\setbeamertemplate{frametitle}{}
\begin{frame}[plain]
  \begin{tikzpicture}[remember picture, overlay]
    % Fondo degradado lateral
    \fill[blooddk, opacity=.12]
      (current page.north west) rectangle (current page.south east);
    % Mancha roja decorativa derecha
    \fill[bloodred, opacity=.10]
      ([xshift=2cm]current page.east) circle (6cm);
    % Mancha roja decorativa izquierda inferior
    \fill[blooddk, opacity=.08]
      ([xshift=-3cm, yshift=-2cm]current page.south west) circle (5cm);
    % Línea lateral roja
    \fill[bloodred]
      ([xshift=0.55\paperwidth]current page.south west)
      rectangle
      ([xshift=0.55\paperwidth+3pt]current page.north west);
  \end{tikzpicture}

  \vspace*{0.8cm}
  \begin{columns}[T]
    \begin{column}{0.52\textwidth}
      \vspace{0.6cm}
      % Tag
      \tagbadge{Aplicación Móvil · Popayán, Colombia}
      \vspace{0.5cm}

      % Título principal
      {\fontsize{42}{44}\selectfont\bfseries\color{bloodlt}!Blood}

      \vspace{4pt}
      {\color{offwhite}\large\itshape ``Dona vida, recibe vida''}

      \vspace{10pt}
      \begin{beamercolorbox}[wd=3cm, ht=2pt]{separation line}\end{beamercolorbox}
      \vspace{10pt}

      {\footnotesize\color{warmgray}
        \textbf{\color{offwhite}Ana María Campo Agredo}\\[3pt]
        Universidad del Cauca\\
        Fac. Ciencias Naturales, Exactas y de la Educación\\
        Dpto. de Física --- Popayán, Colombia\\[3pt]
        \textbf{\color{offwhite}2026}
      }

      \vspace{12pt}
      \tagbadge{\faDroplet\enspace Donación de Sangre}\quad
      \tagbadge{\faMobileAlt\enspace App Móvil}
    \end{column}
    \begin{column}{0.44\textwidth}
      \vspace*{0.4cm}
      \centering
      % Gota grande
      \begin{tikzpicture}
        \fill[top color=bloodlt, bottom color=blooddk,
              drop shadow={opacity=.5, shadow xshift=2pt, shadow yshift=-4pt}]
          (0,3.5) .. controls (-1.9,1.2) and (-2.4,-1) ..
          (0,-2.4) .. controls (2.4,-1) and (1.9,1.2) ..
          (0,3.5);
        \fill[white, opacity=.18]
          (-0.6,0.6) .. controls (-1.2,0.15) and (-1.1,-0.4) ..
          (-0.5,-0.65) .. controls (-0.1,-0.85) and (0.15,-0.6) ..
          (-0.6,0.6);
        \node[white, font=\bfseries\large] at (0,-0.3) {!Blood};
      \end{tikzpicture}
    \end{column}
  \end{columns}
\end{frame}
}

% ──────────────────────────────────────────────────────
%  DIAPOSITIVA 2 – AGENDA
% ──────────────────────────────────────────────────────
\begin{frame}{Agenda}
  \begin{columns}[T]
    \begin{column}{0.48\textwidth}
      \begin{enumerate}[\color{bloodred}1.]\color{offwhite}
        \setlength\itemsep{7pt}
        \item \textbf{Contexto y Crisis} — OMS y Cauca
        \item \textbf{Problemática Identificada} — 6 causas
        \item \textbf{La Solución} — ¿Qué es !Blood?
        \item \textbf{Módulos Funcionales} — ¿Qué hace?
        \item \textbf{Perfiles de Usuario} — ¿Para quién?
        \item \textbf{Prueba de Usabilidad} — Thinking Aloud
      \end{enumerate}
    \end{column}
    \begin{column}{0.48\textwidth}
      \begin{enumerate}[\color{bloodred}1.]\color{offwhite}
        \setlength\itemsep{7pt}
        \setcounter{enumi}{6}
        \item \textbf{Disponibilidad} — Heartbeat + Degradation
        \item \textbf{Modificabilidad} — Dividir + Encapsular
        \item \textbf{Usabilidad} — Undo + Modelos
        \item \textbf{Conclusiones}
      \end{enumerate}
    \end{column}
  \end{columns}
  \vfill
  \begin{tcolorbox}[graybox]
    \small Tiempo estimado de exposición: \redtext{\textbf{10 minutos}} \quad·\quad 10 diapositivas \quad·\quad $\approx$60 seg/diapositiva
  \end{tcolorbox}
\end{frame}

% ──────────────────────────────────────────────────────
%  DIAPOSITIVA 3 – CONTEXTO Y CRISIS
% ──────────────────────────────────────────────────────
\begin{frame}{El Problema}{\tagbadge{Contexto}}
  \begin{columns}[T]
    \begin{column}{0.55\textwidth}
      \begin{tcolorbox}[graybox]
        \small La donación de sangre es un \textbf{pilar fundamental} de los sistemas de salud.
        En Colombia y el Cauca, los niveles están \redtext{significativamente por debajo} del umbral mínimo establecido por la OMS.
      \end{tcolorbox}
      \vspace{8pt}
      \begin{tcolorbox}[graybox]
        \small En Popayán, la \textbf{dispersión geográfica} y la \textbf{ausencia de mecanismos ágiles} impiden conectar donantes con pacientes en situaciones urgentes.
      \end{tcolorbox}
    \end{column}
    \begin{column}{0.42\textwidth}
      % Tres cajas de estadística
      \begin{tcolorbox}[statbox]
        {\color{bloodlt}\fontsize{28}{30}\selectfont\bfseries 10--20}\\[2pt]
        \greytext{\tiny donaciones/1\,000 hab.\\ requeridas por la OMS}
      \end{tcolorbox}
      \vspace{6pt}
      \begin{tcolorbox}[statbox]
        {\color{bloodlt}\fontsize{28}{30}\selectfont\bfseries$\downarrow$}\\[2pt]
        \greytext{\tiny Colombia y Cauca\\ muy por debajo del umbral}
      \end{tcolorbox}
      \vspace{6pt}
      \begin{tcolorbox}[statbox]
        {\color{bloodlt}\fontsize{28}{30}\selectfont\bfseries 0}\\[2pt]
        \greytext{\tiny canales unificados\\ de convocatoria oficial}
      \end{tcolorbox}
    \end{column}
  \end{columns}
\end{frame}

% ──────────────────────────────────────────────────────
%  DIAPOSITIVA 4 – 6 PROBLEMAS
% ──────────────────────────────────────────────────────
\begin{frame}{Problemática Identificada}{\tagbadge{6 causas raíz}}
  \begin{columns}[T]
    \begin{column}{0.48\textwidth}
      \begin{tcolorbox}[redbox, title={\faSignal\enspace Información Dispersa}]
        Sin canal unificado: convocatorias en redes sociales, volantes o comunicados sin fuente oficial.
      \end{tcolorbox}
      \vspace{4pt}
      \begin{tcolorbox}[redbox, title={\faUnlink\enspace Oferta $\neq$ Demanda}]
        Cadenas informales en WhatsApp sin filtro por tipo de sangre ni cercanía geográfica.
      \end{tcolorbox}
      \vspace{4pt}
      \begin{tcolorbox}[redbox, title={\faBrain\enspace Mitos y Desinformación}]
        Creencias falsas disuaden a donantes potenciales dispuestos a donar con información correcta.
      \end{tcolorbox}
    \end{column}
    \begin{column}{0.48\textwidth}
      \begin{tcolorbox}[redbox, title={\faBell\enspace Sin Seguimiento}]
        El donante no recibe retroalimentación, no tiene historial ni recordatorios de recurrencia.
      \end{tcolorbox}
      \vspace{4pt}
      \begin{tcolorbox}[redbox, title={\faHospital\enspace Sin Herramientas Pro}]
        Los centros de salud carecen de medio digital para gestionar campañas y cupos.
      \end{tcolorbox}
      \vspace{4pt}
      \begin{tcolorbox}[redbox, title={\faLock\enspace Comunicación Insegura}]
        Los canales informales no garantizan privacidad ni limitan interacciones a fines de donación.
      \end{tcolorbox}
    \end{column}
  \end{columns}
\end{frame}

% ──────────────────────────────────────────────────────
%  DIAPOSITIVA 5 – LA SOLUCIÓN
% ──────────────────────────────────────────────────────
\begin{frame}{La Solución}{\tagbadge{!Blood – App Móvil}}
  \begin{columns}[T]
    \begin{column}{0.60\textwidth}
      \vspace{4pt}
      \begin{tcolorbox}[graybox]
        \textbf{!Blood} es una aplicación móvil que \textbf{centraliza} solicitudes, campañas y comunicación entre usuarios, priorizando la \redtext{compatibilidad por tipo de sangre} y la \redtext{cercanía geográfica} para facilitar respuestas oportunas.
      \end{tcolorbox}
      \vspace{8pt}
      \begin{block}{\faUsers\enspace 2 Perfiles diferenciados}
        Ciudadano y Profesional de Salud --- funcionalidades propias para cada uno.
      \end{block}
      \vspace{4pt}
      \begin{block}{\faMapMarkerAlt\enspace Conexión inteligente}
        Matching por grupo sanguíneo + geolocalización en tiempo real.
      \end{block}
      \vspace{4pt}
      \begin{block}{\faShieldAlt\enspace Entorno seguro y moderado}
        Verificación de identidad, privacidad de datos y moderación de contenido.
      \end{block}
    \end{column}
    \begin{column}{0.36\textwidth}
      \centering
      \vspace{0.4cm}
      \begin{tikzpicture}
        \fill[top color=bloodlt, bottom color=blooddk,
              drop shadow={opacity=.6, shadow xshift=3pt, shadow yshift=-5pt}]
          (0,3.2) .. controls (-1.75,1.1) and (-2.2,-0.9) ..
          (0,-2.2) .. controls (2.2,-0.9) and (1.75,1.1) ..
          (0,3.2);
        \fill[white, opacity=.2]
          (-0.55,0.55) .. controls (-1.1,0.12) and (-1.0,-0.35) ..
          (-0.45,-0.60) .. controls (-0.1,-0.78) and (0.12,-0.55) ..
          (-0.55,0.55);
        \node[white, font=\bfseries\large] at (0,-0.2) {!Blood};
        \node[white, font=\tiny\itshape] at (0,-0.85) {Dona vida,};
        \node[white, font=\tiny\itshape] at (0,-1.2) {recibe vida};
      \end{tikzpicture}
    \end{column}
  \end{columns}
\end{frame}

% ──────────────────────────────────────────────────────
%  DIAPOSITIVA 6 – MÓDULOS FUNCIONALES
% ──────────────────────────────────────────────────────
\begin{frame}{Módulos Funcionales de la App}{\tagbadge{¿Qué hace !Blood?}}
  \begin{columns}[T]
    \begin{column}{0.48\textwidth}
      \begin{itemize}\setlength\itemsep{5pt}
        \item[\color{bloodred}\faGlobe] \textbf{Splash \& Onboarding} --- primera experiencia
        \item[\color{bloodred}\faKey] \textbf{Login \& Registro} --- 3 pasos + verificación de identidad
        \item[\color{bloodred}\faClipboardList] \textbf{Encuesta de Aptitud} --- filtro automático de requisitos
        \item[\color{bloodred}\faExclamationTriangle] \textbf{Solicitudes Urgentes} --- publicación y notificación inmediata
        \item[\color{bloodred}\faCalendarAlt] \textbf{Campañas de Donación} --- creadas por profesionales de salud
      \end{itemize}
    \end{column}
    \begin{column}{0.48\textwidth}
      \begin{itemize}\setlength\itemsep{5pt}
        \item[\color{bloodred}\faMapMarkerAlt] \textbf{Geolocalización} --- matching por type + proximidad
        \item[\color{bloodred}\faHistory] \textbf{Historial Personal} --- seguimiento + recordatorios automáticos
        \item[\color{bloodred}\faBell] \textbf{Notificaciones Push} --- alertas de compatibilidad cercana
        \item[\color{bloodred}\faTrophy] \textbf{Gamificación} --- insignias, rachas y recompensas
        \item[\color{bloodred}\faFlag] \textbf{Moderación} --- reportes y control de contenido
      \end{itemize}
    \end{column}
  \end{columns}
  \vfill
  \begin{tcolorbox}[graybox]
    \small El registro incluye \textbf{foto en vivo} (anti-duplicados), carga de documento de identidad y geolocalización automática.
  \end{tcolorbox}
\end{frame}

% ──────────────────────────────────────────────────────
%  DIAPOSITIVA 7 – PERFILES DE USUARIO
% ──────────────────────────────────────────────────────
\begin{frame}{Perfiles de Usuario}{\tagbadge{¿Para quién?}}
  \begin{columns}[T]
    % Perfil 1
    \begin{column}{0.31\textwidth}
      \begin{tcolorbox}[
        enhanced, colback=bgcard, colframe=bloodred, arc=10pt,
        coltext=offwhite, fonttitle=\bfseries\small,
        title={\faHeartbeat\enspace Donante Potencial},
        coltitle=offwhite
      ]
        \greytext{\tiny 18–65 años · Popayán · Peso >50 kg}\\[5pt]
        Quiere ayudar donando. Necesita claridad sobre requisitos, dónde donar y recordatorios automáticos.\\[5pt]
        \redtext{\faInfoCircle\enspace} Hombre: cada 3 meses\\
        \redtext{\faInfoCircle\enspace} Mujer: cada 6 meses
      \end{tcolorbox}
    \end{column}
    % Perfil 2
    \begin{column}{0.31\textwidth}
      \begin{tcolorbox}[
        enhanced, colback=bgcard, colframe=accentpink, arc=10pt,
        coltext=offwhite, fonttitle=\bfseries\small,
        title={\faSos\enspace Solicitante de Sangre},
        coltitle=offwhite
      ]
        \greytext{\tiny Familiar/cuidador de paciente}\\[5pt]
        Necesita donantes \textbf{compatibles y cercanos} de urgencia. El proceso debe ser \redtext{extremadamente rápido e intuitivo}.\\[5pt]
        \redtext{\faClock\enspace} Capacidad de atención reducida en emergencia
      \end{tcolorbox}
    \end{column}
    % Perfil 3
    \begin{column}{0.31\textwidth}
      \begin{tcolorbox}[
        enhanced, colback=bgcard, colframe=warmgray, arc=10pt,
        coltext=offwhite, fonttitle=\bfseries\small,
        title={\faHospital\enspace Profesional de Salud},
        coltitle=offwhite
      ]
        \greytext{\tiny Hospital · Clínica · Banco de Sangre}\\[5pt]
        Crea \textbf{campañas institucionales}, gestiona cupos y comunica requisitos de manera estandarizada.\\[5pt]
        \redtext{\faCheckCircle\enspace} Cuenta verificada y contenido oficial
      \end{tcolorbox}
    \end{column}
  \end{columns}
\end{frame}

% ──────────────────────────────────────────────────────
%  DIAPOSITIVA 8 – PRUEBA DE USABILIDAD
% ──────────────────────────────────────────────────────
\begin{frame}{Prueba de Usabilidad}{\tagbadge{Thinking Aloud}}
  \begin{columns}[T]
    \begin{column}{0.32\textwidth}
      \begin{tcolorbox}[statbox]
        {\color{bloodlt}\fontsize{26}{28}\selectfont\bfseries 4}\\[2pt]
        \greytext{\tiny Participantes\\Estudiantes Ing.~Física}
      \end{tcolorbox}
      \vspace{6pt}
      \begin{tcolorbox}[statbox]
        {\color{bloodlt}\fontsize{26}{28}\selectfont\bfseries 12}\\[2pt]
        \greytext{\tiny Hallazgos clave\\identificados}
      \end{tcolorbox}
      \vspace{6pt}
      \begin{tcolorbox}[graybox]
        \tiny Técnica: expresión en voz alta del pensamiento durante la interacción con la interfaz.
      \end{tcolorbox}
    \end{column}
    \begin{column}{0.65\textwidth}
      \footnotesize
      \renewcommand{\arraystretch}{1.25}
      \begin{tabular}{>{\color{bloodlt}\bfseries\footnotesize}p{2cm}
                      >{\color{offwhite}\footnotesize}p{3.4cm}
                      >{\color{warmgray}\footnotesize}p{3.4cm}}
        \rowcolor{blooddk!60}
        \color{offwhite}Categoría & \color{offwhite}Hallazgo & \color{offwhite}Recomendación \\
        \hline
        Registro      & Demasiados datos al inicio & Simplificar el registro \\
        Navegación    & Botón ``Donar'' poco visible & Ubicar en zona destacada \\
        Urgencia      & No es claro cuándo es urgente & Apoyar con evidencia \\
        Motivación    & No mantiene interés a largo plazo & Gamificación \\
        Reversión     & No permite revertir decisión & Permitir con advertencia \\
        Identidad     & Nombres ``creativos'' restán seriedad & Usar nombres reales \\
      \end{tabular}
    \end{column}
  \end{columns}
\end{frame}

% ──────────────────────────────────────────────────────
%  DIAPOSITIVA 9 – DISPONIBILIDAD
% ──────────────────────────────────────────────────────
\begin{frame}{Atributo de Calidad \#1}{\tagbadge{Disponibilidad}}
  \begin{tcolorbox}[graybox]
    \small\textbf{Caso crítico:} El servidor de notificaciones deja de responder durante una \redtext{solicitud urgente de donación}.
  \end{tcolorbox}
  \vspace{6pt}
  \begin{columns}[T]
    \begin{column}{0.48\textwidth}
      \begin{block}{\faHeartbeat\enspace Táctica 1: Heartbeat}
        Señales periódicas entre monitor y proceso supervisado. Detecta fallos en segundos en:
        \begin{itemize}\setlength\itemsep{2pt}
          \item Servicio de notificaciones urgentes
          \item Módulo de mensajería
          \item Backend de solicitudes de sangre
          \item Conexión con base de datos
        \end{itemize}
      \end{block}
    \end{column}
    \begin{column}{0.48\textwidth}
      \begin{block}{\faArrowDown\enspace Táctica 2: Degradation}
        Ante fallo, mantiene funciones críticas y suspende secundarias:
        \begin{itemize}\setlength\itemsep{2pt}
          \item Falla chat \faArrowRight{} solicitudes siguen activas
          \item Falla notif. \faArrowRight{} campañas visibles en app
          \item Priorizar: registro, autenticación, compatibilidad
          \item Reportes estadísticos pueden fallar sin impacto
        \end{itemize}
      \end{block}
    \end{column}
  \end{columns}
\end{frame}

% ──────────────────────────────────────────────────────
%  DIAPOSITIVA 10 – MODIFICABILIDAD
% ──────────────────────────────────────────────────────
\begin{frame}{Atributo de Calidad \#2}{\tagbadge{Modificabilidad}}
  \begin{tcolorbox}[graybox]
    \small\textbf{Caso:} Se detecta uso inapropiado del \textbf{chat en vivo} — usuarios comparten teléfonos o promueven compra/venta de sangre. El equipo debe \redtext{eliminar el módulo sin afectar el resto del sistema}.
  \end{tcolorbox}
  \vspace{6pt}
  \begin{columns}[T]
    \begin{column}{0.48\textwidth}
      \begin{block}{\faCut\enspace Táctica 1: Dividir Módulo}
        Cada funcionalidad tiene responsabilidad única e independiente:
        \begin{itemize}\setlength\itemsep{2pt}
          \item Chat $\neq$ Solicitudes de donación
          \item Registro $\neq$ Gestión de perfiles
          \item Campañas $\neq$ Historial de donaciones
          \item Notificaciones $\neq$ Búsqueda de donantes
        \end{itemize}
        \redtext{Resultado:} eliminar chat no rompe ninguna otra función.
      \end{block}
    \end{column}
    \begin{column}{0.48\textwidth}
      \begin{block}{\faBox\enspace Táctica 2: Encapsular}
        Cada módulo expone solo su API, ocultando implementación interna:
        \begin{itemize}\setlength\itemsep{2pt}
          \item Chat: enviar, bloquear, cerrar conversación
          \item Notificaciones: enviar alerta, marcar leída
          \item Compatibilidad: validar grupo sanguíneo
          \item Campañas: publicar, actualizar, cerrar jornada
        \end{itemize}
        \redtext{Resultado:} cambios internos no se propagan.
      \end{block}
    \end{column}
  \end{columns}
\end{frame}

% ──────────────────────────────────────────────────────
%  DIAPOSITIVA 11 – USABILIDAD
% ──────────────────────────────────────────────────────
\begin{frame}{Atributo de Calidad \#3}{\tagbadge{Usabilidad}}
  \begin{tcolorbox}[graybox]
    \small\textbf{Caso:} Un donante confirma su intención de donar y luego desea \redtext{revertir esa decisión} por un error o cambio de disponibilidad.
  \end{tcolorbox}
  \vspace{6pt}
  \begin{columns}[T]
    \begin{column}{0.31\textwidth}
      \begin{block}{\faUndo\enspace Undo (Deshacer)}
        El sistema conserva el estado anterior. El donante puede revertir \textbf{con advertencia}, sin inconsistencias. Restricción parcial para no afectar coordinación médica.
      \end{block}
    \end{column}
    \begin{column}{0.31\textwidth}
      \begin{block}{\faChartBar\enspace Modelo del Sistema}
        Retroalimentación constante:
        \begin{itemize}\setlength\itemsep{2pt}
          \item ``23 donantes notificados''
          \item ``2 de 4 confirmados''
          \item Progreso en registro (3 pasos)
          \item Cupos disponibles en campañas
        \end{itemize}
      \end{block}
    \end{column}
    \begin{column}{0.31\textwidth}
      \begin{block}{\faPuzzlePiece\enspace Modelo de Tarea}
        La app adapta la interfaz al contexto:
        \begin{itemize}\setlength\itemsep{2pt}
          \item Donante \faArrowRight{} checklist salud
          \item Solicitante \faArrowRight{} campos urgentes primero
          \item Profesional \faArrowRight{} cupos + horario
        \end{itemize}
      \end{block}
    \end{column}
  \end{columns}
\end{frame}

% ──────────────────────────────────────────────────────
%  DIAPOSITIVA 12 – CONCLUSIONES
% ──────────────────────────────────────────────────────
{
\setbeamertemplate{footline}{}
\begin{frame}[plain]{Conclusiones}
  \begin{tikzpicture}[remember picture, overlay]
    \fill[bloodred, opacity=.10]
      ([xshift=3cm]current page.east) circle (7cm);
    \fill[blooddk, opacity=.08]
      ([xshift=-3cm, yshift=-3cm]current page.south west) circle (5cm);
  \end{tikzpicture}

  \vspace{6pt}
  \begin{columns}[T]
    \begin{column}{0.60\textwidth}
      {\large\bfseries\color{bloodlt}!Blood}\greytext{ — Tecnología que Salva Vidas}
      \vspace{8pt}

      \begin{itemize}\setlength\itemsep{7pt}
        \item[\color{bloodred}\faDroplet]
          \textbf{Centraliza} la donación de sangre en Popayán y el Cauca
        \item[\color{bloodred}\faMapMarkerAlt]
          Matching \textbf{inteligente} por tipo de sangre y proximidad
        \item[\color{bloodred}\faShieldAlt]
          Entorno \textbf{seguro, moderado y verificado}
        \item[\color{bloodred}\faRecycle]
          Sistema \textbf{resiliente} (Disponibilidad) y \textbf{adaptable} (Modificabilidad)
        \item[\color{bloodred}\faStar]
          Diseñado desde la \textbf{usabilidad real} validada con usuarios
      \end{itemize}

      \vspace{10pt}
      \begin{tcolorbox}[graybox]
        Una solución móvil que \textbf{reduce barreras de acceso}, incrementa la donación voluntaria y contribuye a \redtext{\textbf{salvar vidas}}.
      \end{tcolorbox}
    \end{column}
    \begin{column}{0.36\textwidth}
      \centering
      \vspace{0.6cm}
      \begin{tikzpicture}
        \fill[top color=bloodlt, bottom color=blooddk,
              drop shadow={opacity=.7, shadow xshift=4pt, shadow yshift=-6pt}]
          (0,3.5) .. controls (-1.9,1.2) and (-2.4,-1.0) ..
          (0,-2.4) .. controls (2.4,-1.0) and (1.9,1.2) ..
          (0,3.5);
        \fill[white, opacity=.2]
          (-0.6,0.6) .. controls (-1.2,0.15) and (-1.1,-0.4) ..
          (-0.5,-0.65) .. controls (-0.1,-0.85) and (0.15,-0.6) ..
          (-0.6,0.6);
        \node[white, font=\bfseries\Large] at (0,-0.1) {!Blood};
        \node[white, font=\scriptsize\itshape] at (0,-0.8) {``Dona vida,};
        \node[white, font=\scriptsize\itshape] at (0,-1.2) {recibe vida''};
      \end{tikzpicture}

      \vspace{8pt}
      {\small\color{warmgray}\itshape Ana María Campo Agredo\\
      Universidad del Cauca · 2026}
    \end{column}
  \end{columns}
\end{frame}
}

% ════════════════════════════════════════════════════════
\end{document}
% ════════════════════════════════════════════════════════
